\begin{afterword}
\summaryhead

The author tells of meeting a man who called himself purely selfish and urged the author to be selfish too. The author answered that religious people seem to find religious reasons for whatever they already want to do, and suggested that a philosophy of selfishness works the same way: it has no effect on behaviour, because its holders can rationalize their kindness in selfish terms.

When the man disagreed, the author asked why a selfish man would want others to be selfish. He replied that selfish people would see that their interest lies in productive work, not in laws against his property. The author answered that the author is already a libertarian and has taken a lower-paid job to benefit others, asked whether converting people is the most selfish use of the man's time, and asked whether his wish to convert came before his egoism. The man conceded that it may have. A footnote offers three more test questions for professed egoists.

\responsehead

The thesis is offered as an impression. ``It looks to me like'' a philosophy of selfishness ``has no effect on their behavior,'' and the comparison with religion is based on the believers the author has met. The evidence given is one conversation. The conversation bears on a narrower point than the thesis. It concerns one act, urging selfishness on the author, and the man's reasons for it.

What the conversation does show comes from the man's concession: his wish to convert others may have come first, and his selfish reason for it second. That supports the claim that his argument was a rationalization.

The man's answer does answer the author's first question. His reason for wanting others to be selfish is itself self-interested, and it is a familiar one. Ayn Rand held that the rational interests of different people do not conflict; on that view, wanting others to be rationally selfish fits egoism. The author questions whether it applies in the author's own case, which is a reasonable reply to a man speaking to the author. The final question, about which aim came first, is the one that tests the claim, and the man grants that the author ``may be right'' about it. The first question, why an egoist would recommend egoism, echoes a point the author had made in ``Your Rationality Is My Business''; see the notes there.

The footnote's questions are well made. Each includes a clause that rules out the usual reply. The stipulation that the egoist will not know about the torture rules out the reply that helping others only relieves one's own distress, a reply that Batson's experiments tested and did not support. The undetectable theft is an old test, the ring of Gyges in Plato's \textsc{Republic}.

In short: a hedged suspicion about professed egoists, supported by one conversation that shows one man rationalizing his wish to convert others, and followed by well-made test questions, one of them as old as Plato.
\end{afterword}
